% =====================================================================
%  TCET CODING PLATFORM -- Official Technical Documentation
%  MODULE 3 of 7 : Authentication, Authorization, Security & Data Persistence
%  STANDALONE FILE.  Compile: pdflatex Module3_Security_Data.tex  (twice)
% =====================================================================
\documentclass[12pt,a4paper,oneside]{report}
\usepackage[T1]{fontenc}
\usepackage[utf8]{inputenc}
\usepackage[a4paper,margin=1in]{geometry}
\usepackage{mathptmx}\usepackage{microtype}\usepackage{graphicx}\usepackage{xcolor}
\usepackage{tabularx}\usepackage{booktabs}\usepackage{longtable}\usepackage{xltabular}
\usepackage{array}\usepackage{enumitem}\usepackage{listings}\usepackage{titlesec}
\usepackage{fancyhdr}\usepackage{parskip}\usepackage[hidelinks]{hyperref}
\definecolor{tcetblue}{RGB}{0,0,0}\definecolor{tcetaccent}{RGB}{0,0,0}
\definecolor{codebg}{RGB}{245,246,248}\definecolor{codeframe}{RGB}{215,219,225}
\definecolor{codekw}{RGB}{0,0,0}\definecolor{codecomment}{RGB}{0,0,0}
\definecolor{codestr}{RGB}{0,0,0}\definecolor{boxframe}{RGB}{0,0,0}
\hypersetup{colorlinks=true,linkcolor=tcetblue,urlcolor=tcetaccent,citecolor=tcetblue,
  pdftitle={TCET Coding Platform -- Module 3: Security \& Data},pdfauthor={Centre of Excellence, TCET}}
\lstdefinestyle{tcp}{backgroundcolor=\color{codebg},basicstyle=\ttfamily\footnotesize,
  keywordstyle=\color{codekw}\bfseries,commentstyle=\color{codecomment}\itshape,
  stringstyle=\color{codestr},numberstyle=\tiny\color{gray},frame=single,
  rulecolor=\color{codeframe},breaklines=true,showstringspaces=false,tabsize=2,
  columns=fullflexible,keepspaces=true,xleftmargin=0.6em,framexleftmargin=0.6em}
\lstset{style=tcp}
\newcommand{\code}[1]{\texttt{#1}}
\newenvironment{callout}[1]{\par\vspace{4pt}\noindent
  \begin{tabular}{@{}p{0.014\linewidth}@{\hspace{0.6em}}p{0.93\linewidth}@{}}
  \textcolor{boxframe}{\rule{2pt}{\baselineskip}} & \textbf{\textcolor{tcetblue}{#1}}\\[2pt] & }{\end{tabular}\par\vspace{6pt}}
\titleformat{\chapter}[display]{\bfseries\fontsize{18}{22}\selectfont}{\bfseries\fontsize{14}{17}\selectfont Chapter \thechapter}{0.2em}{}
\titlespacing*{\chapter}{0pt}{6pt}{22pt}
\titleformat{\section}{\bfseries\fontsize{16}{19}\selectfont}{\thesection}{0.6em}{}
\titleformat{\subsection}{\bfseries\fontsize{14}{17}\selectfont}{\thesubsection}{0.5em}{}
\pagestyle{fancy}\fancyhf{}\renewcommand{\headrulewidth}{0.4pt}
\fancyhead[L]{\small\itshape TCET Coding Platform}\fancyhead[R]{\small\itshape Module 3 \textbullet\ Security \& Data}
\fancyfoot[C]{\thepage}
\fancypagestyle{plain}{\fancyhf{}\fancyfoot[C]{\thepage}\renewcommand{\headrulewidth}{0pt}}
\newcommand{\note}[1]{\par\vspace{3pt}\noindent{\color{tcetaccent}\textbf{Note.}}~#1\par\vspace{3pt}}
\begin{document}
\begin{titlepage}\centering
\vspace*{1.4cm}{\color{tcetaccent}\rule{\linewidth}{1.2pt}}\\[0.5cm]
{\LARGE\bfseries TCET Coding Platform\par}
\vspace{0.4cm}{\Large Official Technical Documentation\par}{\large Draft~1 \textemdash\ Testing Phase\par}
{\color{tcetaccent}\rule{\linewidth}{1.2pt}}\\[1.4cm]
{\Huge\bfseries\color{tcetblue} Module 3\par}\vspace{0.3cm}
{\Huge Authentication, Authorization,\\[0.2cm] Security \& Data Persistence\par}\vspace{2.0cm}
{\large\bfseries Powered by\par}{\Large\color{tcetblue} Centre of Excellence (CoE), TCET\par}\vfill
\begin{tabular}{rl}\textbf{Module} & 3 of 7\\[3pt]
\textbf{Scope} & Identity, RBAC, hardening, data model\\[3pt]
\textbf{Status} & Draft~1 (living document until Stage~2)\\[3pt]
\textbf{Audience} & Security reviewers, Officials, Engineering\\\end{tabular}
\vspace{0.8cm}{\small\itshape \today\par}\end{titlepage}
\pagenumbering{roman}\tableofcontents\clearpage\pagenumbering{arabic}

% =====================================================================
\chapter{Overview \& Security Objectives}
% =====================================================================
\section{Scope of this Module}
This module specifies how the platform establishes \emph{who} a caller is
(authentication), what they are permitted to do (authorization), how the attack
surface is hardened (security controls), and how all state is persisted (the data
layer). Together these constitute the platform's confidentiality, integrity, and
availability posture at the identity and data tiers.

\section{Security Objectives}
\begin{longtable}{@{}p{0.28\linewidth} p{0.64\linewidth}@{}}
\toprule
\textbf{Objective} & \textbf{Statement}\\
\midrule
\endhead
Authenticated access only & No request reaches domain logic without a verified, active identity forwarded by a trusted proxy.\\
No credential custody & The platform stores no passwords; identity is federated from the CoE SSO.\\
Least privilege & Each role can perform only what its function requires, enforced at multiple layers.\\
Tamper resistance & Scoring, violation counts, and results are computed server-side from persisted state.\\
Data integrity & Persisted documents conform to a validated schema; history is stable via snapshots.\\
Abuse resistance & Rate limiting, payload caps, and origin checks blunt floods and cross-site abuse.\\
\bottomrule
\end{longtable}

\section{Security Design Principles}
\begin{description}[leftmargin=1.6em,style=nextline]
  \item[Defence in depth] Independent controls at the proxy, identity,
        authorization, and data layers.
  \item[Fail closed] Ambiguous or unverifiable requests are rejected, not
        admitted.
  \item[Server-side authority] The client is never trusted for anything that
        affects scoring, eligibility, or results.
  \item[Minimal exposure] Internal and diagnostic endpoints are restricted to
        trusted sources; sensitive data never appears in URLs.
\end{description}

% =====================================================================
\chapter{Identity \& the Trusted-Proxy Model}
% =====================================================================
\section{No Local Credentials}
The platform does not manage passwords or credentials. Authentication is
performed upstream by the CoE centralized SSO. The backend consumes a
\emph{verified} identity rather than verifying credentials itself.

\section{The Trust Boundary}
The architecture is a \textbf{trusted reverse-proxy model}:
\begin{enumerate}[leftmargin=1.6em]
  \item The user authenticates against the CoE SSO at the edge.
  \item The proxy forwards only authenticated requests to the backend, attaching
        a signed identity as headers and/or a token cookie.
  \item The backend verifies that the request arrived from a trusted proxy source
        and then trusts the forwarded identity.
\end{enumerate}

\begin{callout}{Why the source-IP check is load-bearing}
Header-based identity is safe \emph{only} because the backend independently
verifies that the request came from an allowlisted proxy source address. If the
backend were exposed directly to the internet, a client could fabricate identity
headers. The source check is therefore the single control that makes the whole
model sound, and it runs on every request.
\end{callout}

\section{Identity Signals}
\begin{longtable}{@{}p{0.26\linewidth} p{0.66\linewidth}@{}}
\toprule
\textbf{Signal} & \textbf{Description}\\
\midrule
\endhead
Identity headers & \code{x-coe-email}, \code{x-coe-name}, \code{x-coe-role}, \code{x-coe-status} injected by the trusted proxy.\\
Signed token (fallback) & A CoE JWT taken from a dedicated header, a bearer authorization header, or one of several federated cookies, verified with a strong algorithm against a configured secret.\\
\bottomrule
\end{longtable}

When the required headers are present they are used directly; otherwise the
middleware verifies the signed token. Either way, the resolved identity must
carry a valid email, a recognized role, and an active status.

\section{Authentication Decision Flow}
\begin{lstlisting}[language={}]
1. Is the request source IP in the trusted-proxy allowlist?
      no  -> 401 "Unauthorized source"  (+ security log)
2. Are all identity headers present?
      yes -> validate email format and role value
      no  -> verify the signed token; reject if invalid -> 401
3. Is the resolved role recognized?
      no  -> 401 "invalid authentication token"
4. Is the account status ACTIVE?
      no  -> 403 "Account is <STATUS>"
5. Synchronize the user record; attach the request user
6. -> continue to the route
\end{lstlisting}

\section{Proxy Source Verification}
The allowlist supports both explicit addresses and CIDR subnets, for IPv4 and
IPv6, and always trusts loopback so local development and internal health checks
work. Invalid entries are ignored with a warning rather than failing the whole
service. Rejections are logged with a masked email and forwarded-chain
diagnostics to aid investigation without exposing personal data.

\section{Session Termination}
Logout clears the federated cookies directly (rather than depending on an
upstream logout page), across the host-only and domain-wide cookie variants, and
redirects to a validated, allowlisted frontend origin.

% =====================================================================
\chapter{Authorization \& Access Control}
% =====================================================================
\section{Role Model}
The CoE issues four roles, which the platform collapses into two internal roles.

\begin{longtable}{@{}p{0.20\linewidth} p{0.22\linewidth} p{0.48\linewidth}@{}}
\toprule
\textbf{CoE role} & \textbf{Platform role} & \textbf{Meaning}\\
\midrule
\endhead
Student & Student & Solves problems; attempts contests.\\
Faculty & Faculty & Authors content; reviews; exports.\\
Admin & Faculty & Treated as faculty in the platform.\\
Industry & Faculty & Treated as faculty in the platform.\\
\bottomrule
\end{longtable}

\section{Layers of Authorization}
\begin{description}[leftmargin=1.6em,style=nextline]
  \item[Route-level role guards] Faculty-only endpoints (problem management,
        standings export, attempt review) are gated by a role requirement.
  \item[Profile-completion guard] Certain actions require a completed profile.
  \item[Ownership \& visibility checks in services] Beyond the role, services
        verify that the specific record belongs to, or is visible to, the caller:
        faculty may only view submissions for resources they own; a student may
        only see published problems for their department; contest access is
        filtered by department.
\end{description}

\begin{callout}{Preventing horizontal escalation}
A passing role check is never sufficient by itself. Even a valid faculty user is
re-checked at the service layer for ownership of the specific resource, so one
faculty member cannot read another's submissions by guessing an identifier, and a
student cannot read another student's data.
\end{callout}

\section{Authorization Matrix (selected endpoints)}
\begin{longtable}{@{}p{0.46\linewidth} p{0.44\linewidth}@{}}
\toprule
\textbf{Action} & \textbf{Permitted to}\\
\midrule
\endhead
List / view published problems & Any authenticated user (student view filtered by department).\\
Create / edit / archive problems & Faculty (owner).\\
Run / submit code & Student.\\
View a submission & The submitting student, or the faculty owner of the resource.\\
Create / edit contests & Faculty (owner).\\
Start / submit a contest attempt & Student (registered).\\
View standings (student) & Student, after publish and feedback.\\
Publish results, review attempts, export CSV & Faculty (owner).\\
View another user's analytics & Faculty.\\
\bottomrule
\end{longtable}

% =====================================================================
\chapter{Security Controls}
% =====================================================================
\section{Control Catalogue}
\begin{longtable}{@{}p{0.28\linewidth} p{0.64\linewidth}@{}}
\toprule
\textbf{Control} & \textbf{Behaviour}\\
\midrule
\endhead
Secure headers & A hardening middleware sets safe response headers; the framework banner is disabled.\\
CORS allowlist & Only configured origins are admitted; credentials are enabled for those origins only.\\
Mutation-origin guard & State-changing requests must carry an allowlisted Origin/Referer, mitigating cross-site request forgery.\\
Frontend-pathname guard & The optional client pathname header is validated against a safe pattern.\\
Payload cap & JSON bodies are limited to 100\,kb.\\
Global rate limiting & Requests are throttled per identity or IP; excess returns HTTP~429.\\
Submission rate limiting & Final submissions and code runs are additionally rate limited per user.\\
Security event logging & Authentication rejections are logged with masked identifiers and forwarded-chain diagnostics.\\
Guarded internal routes & Health and diagnostic routes answer only trusted internal sources; the database diagnostic is disabled in production.\\
\bottomrule
\end{longtable}

\section{Threats \& Mitigations}
\begin{longtable}{@{}p{0.34\linewidth} p{0.58\linewidth}@{}}
\toprule
\textbf{Threat} & \textbf{Mitigation}\\
\midrule
\endhead
Identity spoofing & Trusted-proxy source check; proxy strips client identity headers.\\
Cross-site request forgery & Origin/Referer allowlist on mutations; credentialed CORS restricted to known origins.\\
Request floods / abuse & Global and per-user rate limiting; payload cap.\\
Untrusted code execution & Sandboxed judging with no network and strict limits (Module~4).\\
Result tampering & Server-side scoring and violation accounting from persisted state.\\
Information disclosure & Masked logs; no sensitive data in URLs; restricted internal endpoints.\\
\bottomrule
\end{longtable}

\section{Secrets \& Sensitive Material}
The token-signing secret must meet a minimum strength and is validated at
startup. Environment files and any service-account keys must never be committed
to version control. Personal data is kept out of URLs and query strings, and
emails are masked in security logs.

% =====================================================================
\chapter{Data Architecture}
% =====================================================================
\section{Datastore}
All application state is persisted in MongoDB via the official driver. A single,
lazily-initialized client and database handle is shared process-wide; the
database name is configurable. The store uses \textbf{server-side JSON Schema
validation}: each collection declares required fields and per-field types, so
non-conforming documents are rejected \textemdash\ a second line of defence
behind the application's validation.

\note{Some repository classes carry a legacy \code{Firestore}-prefixed name. The
authoritative and only datastore is MongoDB; the prefix does not indicate a
Firestore dependency.}

\section{Persistence Patterns}
\begin{description}[leftmargin=1.6em,style=nextline]
  \item[Repository per domain] Each domain's repository is the sole gateway to
        its collection(s); services depend on interfaces, not the driver.
  \item[Stable application identifiers] Records carry an application-level
        identifier in addition to the database identifier; writes are idempotent
        upserts keyed on it.
  \item[Denormalized snapshots] Records embed the values they depend on at
        creation time (e.g. a submission stores the problem's title and
        difficulty), keeping history stable and reads cheap.
  \item[Recomputed aggregates] User and problem statistics are recomputed from
        source submissions on finalization, keeping analytics current without a
        batch job.
\end{description}

\section{Collection Map}
\begin{longtable}{@{}p{0.28\linewidth} p{0.64\linewidth}@{}}
\toprule
\textbf{Collection} & \textbf{Contents}\\
\midrule
\endhead
\code{users} & Identity, profile, and denormalized statistics.\\
\code{problems} & Problem bank with tests, limits, lifecycle, and aggregates.\\
\code{submissions} & Every code submission (practice or contest) with verdict and snapshots.\\
\code{leaderboard} & Denormalized ranking projection of user statistics.\\
\code{contests} & Contest definitions with the embedded question set.\\
\code{contest\_registrations} & Per-student contest enrolments.\\
\code{contest\_attempts} & Per-student attempts embedding per-question state.\\
\code{proctoring\_events} & Immutable log of integrity events during attempts.\\
\code{contest\_feedback} & One post-contest questionnaire per student per contest; gates results.\\
\bottomrule
\end{longtable}

% =====================================================================
\chapter{Data Dictionary}
% =====================================================================
This chapter documents the persisted schema of each collection, derived from the
database's own schema definition. Types are given in general terms; ``opt.''
marks an optional field.

\section{Collection: \code{users}}
\begin{xltabular}{\linewidth}{@{}l l X@{}}
\toprule
\textbf{Field} & \textbf{Type} & \textbf{Meaning}\\
\midrule
\endhead
\code{email} & string & Primary identity (lower-cased).\\
\code{role} & string & Student or Faculty.\\
\code{name} & string & Display name.\\
\code{uid} & string / null & Institutional user id.\\
\code{isProfileComplete} & bool & Whether onboarding is finished.\\
\code{designation} & string / null & Faculty designation.\\
\code{rollNumber} & string / null & Student roll number.\\
\code{department} & string / null & Academic department.\\
\code{semester} & int / null & Current semester (drives ``year'').\\
\code{linkedInUrl}, \code{githubUrl} & string / null & Profile links.\\
\code{skills} & array & Declared skills.\\
\code{rating}, \code{score} & int & Difficulty-weighted rating (score mirrors rating).\\
\code{problemsSolved} & int & Distinct problems accepted.\\
\code{submissionCount} & int & Finalized problem submissions.\\
\code{acceptedSubmissionCount} & int & Accepted submissions.\\
\code{accuracy} & number & Accepted / total, as a percentage.\\
\code{lastAcceptedAt} & date / null & Timestamp of last accepted solve.\\
\code{lastLoginAt} & date & Last authenticated access.\\
\code{createdAt}, \code{updatedAt} & date & Lifecycle timestamps.\\
\bottomrule
\end{xltabular}

\section{Collection: \code{problems}}
\begin{xltabular}{\linewidth}{@{}l l X@{}}
\toprule
\textbf{Field} & \textbf{Type} & \textbf{Meaning}\\
\midrule
\endhead
\code{id}, \code{slug} & string & Application id and URL-friendly identifier.\\
\code{title}, \code{topic} & string & Title and primary topic.\\
\code{statement} & string & Full problem statement.\\
\code{inputFormat}, \code{outputFormat} & string & I/O format descriptions.\\
\code{constraints} & array & Constraint lines.\\
\code{explanation} & string & Worked explanation.\\
\code{examples} & array & Displayed examples (input, output, explanation, hidden).\\
\code{difficulty} & string & Easy / Medium / Hard.\\
\code{tags} & array & Topical tags.\\
\code{timeLimit(Seconds)} & int & Per-run CPU time limit.\\
\code{memoryLimit(Mb)} & int & Per-run memory limit.\\
\code{lifecycleState} & string & Draft / Published / Archived.\\
\code{approved} & bool & Editorial approval flag.\\
\code{targetDepartment} & string / null & Department targeting (null = open).\\
\code{createdBy}, \code{createdByRole} & string & Authoring faculty and role.\\
\code{totalSubmissions}, \code{acceptedSubmissions} & int & Finalized and accepted counts.\\
\code{acceptanceRate} & number & Acceptance percentage.\\
\code{sampleTestCases} & array & Visible tests.\\
\code{hiddenTestCases} & array & Hidden judging tests.\\
\code{createdAt}, \code{updatedAt} & date & Lifecycle timestamps.\\
\bottomrule
\end{xltabular}

\section{Collection: \code{submissions}}
\begin{xltabular}{\linewidth}{@{}l l X@{}}
\toprule
\textbf{Field} & \textbf{Type} & \textbf{Meaning}\\
\midrule
\endhead
\code{id} & string & Application id.\\
\code{sourceType} & string & Practice problem or contest coding.\\
\code{userEmail}, \code{userRole}, \code{userDepartment} & string & Submitter snapshot.\\
\code{resourceOwnerEmail} & string & Owning faculty (visibility control).\\
\code{resourceTargetDepartment} & string / null & Target department snapshot.\\
\code{problemId} & string & Problem or contest-question id.\\
\code{problemTitleSnapshot} & string & Title at submission time.\\
\code{problemDifficultySnapshot} & string & Difficulty at submission time.\\
\code{contestId}, \code{contestTitleSnapshot}, \code{contestQuestionId} & string / null & Contest context, if any.\\
\code{code} & string & Submitted source.\\
\code{language} & string & Submission language.\\
\code{status} & string & Verdict.\\
\code{runtimeMs}, \code{memoryKb} & int & Measured runtime and memory.\\
\code{passedCount}, \code{totalCount} & int & Tests passed / total.\\
\code{executionProvider} & string & Judge0 or stub.\\
\code{ratingAwarded} & int & Rating granted (first accept only).\\
\code{stdout}, \code{stderr} & string / null & Captured output / diagnostics.\\
\code{queueJobId} & string / null & Queue job id.\\
\code{createdAt}, \code{updatedAt} & date & Lifecycle timestamps.\\
\code{judgedAt} & date & When judging completed.\\
\code{finalizationAppliedAt} & date & Idempotency guard for finalization.\\
\bottomrule
\end{xltabular}

\section{Collection: \code{contests}}
\begin{xltabular}{\linewidth}{@{}l l X@{}}
\toprule
\textbf{Field} & \textbf{Type} & \textbf{Meaning}\\
\midrule
\endhead
\code{id}, \code{title} & string & Application id and title.\\
\code{type} & string & Rated or Practice.\\
\code{lifecycleState} & string & Published.\\
\code{startAt}, \code{endAt} & date & Contest window.\\
\code{durationMinutes} & int & Per-attempt duration.\\
\code{registrationOpenAt}, \code{registrationCloseAt} & date & Registration window.\\
\code{resultsPublished} & bool & Whether results are revealed.\\
\code{targetDepartment} & string & Eligibility department.\\
\code{maxViolations} & int & Auto-submit violation threshold.\\
\code{createdBy}, \code{createdByRole} & string & Authoring faculty and role.\\
\code{questions} & array & Embedded question set (see below).\\
\code{createdAt}, \code{updatedAt} & date & Lifecycle timestamps.\\
\bottomrule
\end{xltabular}

\noindent Each embedded \emph{coding} question carries an id, type, points, title,
difficulty, statement, constraints, input/output formats, time and memory limits,
sample and hidden test cases, and supported languages. Objective (MCQ/MSQ)
questions carry a statement, options, and the correct answer(s).

\section{Collection: \code{contest\_attempts}}
\begin{xltabular}{\linewidth}{@{}l l X@{}}
\toprule
\textbf{Field} & \textbf{Type} & \textbf{Meaning}\\
\midrule
\endhead
\code{id}, \code{contestId} & string & Attempt and contest ids.\\
\code{contestTitleSnapshot} & string & Contest title at start.\\
\code{userEmail}, \code{userName}, \code{userUid}, \code{userDepartment} & string & Student snapshot.\\
\code{status} & string & Active / Submitted / Auto-submitted / \ldots\\
\code{score} & int & Awarded points minus penalties.\\
\code{violationCount}, \code{violationPenaltyPoints} & int & Violations and deductions.\\
\code{timeTakenMs} & int & Elapsed time to terminal state.\\
\code{startedAt}, \code{deadlineAt} & date & Start and frozen personal deadline.\\
\code{submittedAt}, \code{autoSubmittedAt}, \code{lastSolvedAt} & date / null & Terminal and solve timestamps.\\
\code{updatedAt} & date & Last mutation.\\
\code{questionStates} & array & Per-question progress (below).\\
\bottomrule
\end{xltabular}

\noindent Each \code{questionStates} entry holds the question id and type, a
status, attempts count, awarded points, the submitted answer or draft code, the
last submission pointer, pass counts, final language/status/runtime/memory, and a
solved timestamp.

\section{Collections: \code{contest\_registrations}, \code{proctoring\_events}, \code{contest\_feedback}}
\begin{xltabular}{\linewidth}{@{}l l X@{}}
\toprule
\textbf{Field} & \textbf{Type} & \textbf{Meaning}\\
\midrule
\endhead
\code{registrations.*} & mixed & Registration id, contest id, student snapshot, and registration time.\\
\code{proctoring.type} & string & Event type (tab switch, fullscreen exit, screenshot, \ldots).\\
\code{proctoring.details} & string & Human-readable detail.\\
\code{proctoring.*} & mixed & Event id, contest/attempt ids, student, and timestamp.\\
\code{feedback.navigationEase} & int (1--5) & Navigation ease rating.\\
\code{feedback.visualDesignRating} & int (1--5) & Visual design rating.\\
\code{feedback.interfaceReadability} & enum & Yes / No / Need improvement.\\
\code{feedback.editorResponsiveness} & int (1--5) & Editor responsiveness.\\
\code{feedback.compilationLag} & int (1--5) & Perceived compile latency.\\
\code{feedback.errorMessageClarity} & int (1--5) & Error-message clarity.\\
\code{feedback.problemStatementClarity} & enum & Yes / No / Needs improvement.\\
\code{feedback.bugsOrBrokenLinks} & string & Free-text bug report.\\
\code{feedback.oneNewFeature} & string & One requested feature.\\
\code{feedback.recommendLikelihood} & int (1--5) & Likelihood to recommend.\\
\code{feedback.overallRating} & int (1--5) / null & Optional overall rating.\\
\bottomrule
\end{xltabular}

% =====================================================================
\chapter{Data Integrity, Privacy \& Lifecycle}
% =====================================================================
\section{Integrity Mechanisms}
\begin{itemize}[leftmargin=1.5em]
  \item \textbf{Schema validation} rejects malformed documents at the database.
  \item \textbf{Idempotent finalization} ensures a verdict is applied once and
        aggregates stay consistent under retries.
  \item \textbf{Server-side scoring} guarantees results reflect persisted state,
        not client claims.
\end{itemize}

\section{Privacy Handling}
Personal data is confined to what the platform needs (identity, profile,
academic details) and is never placed in URLs or query strings. Logs mask email
addresses. Access to another user's data is restricted to faculty and to
resources they own. Data-retention and consent documentation is identified as
future work in Module~7.

\section{Backup \& Retention (current \& planned)}
The current build relies on the operational database's own durability. An
explicit indexing, backup/restore, and archival strategy for large historical
datasets is recorded as future work in Module~7.

% =====================================================================
\chapter{Security Summary}
% =====================================================================
\section{Assurance Summary}
\begin{longtable}{@{}p{0.34\linewidth} p{0.58\linewidth}@{}}
\toprule
\textbf{Property} & \textbf{Basis of assurance}\\
\midrule
\endhead
Authentication & Federated SSO; trusted-proxy source verification; active-status check.\\
Authorization & Role guards plus service-layer ownership/visibility checks.\\
Confidentiality & Restricted internal endpoints; masked logs; no data in URLs.\\
Integrity & Server-side scoring; schema validation; idempotent finalization.\\
Availability & Rate limiting; payload caps; resilience mechanisms (Modules~2, 4, 5).\\
\bottomrule
\end{longtable}

\section{Residual Risks}
Browser-based proctoring cannot control external cameras or a second device;
these are accepted risks addressed by invigilation and future integrations
(Modules~6 and~7). Observability is currently console-based; structured logging
and alerting are future work.

\vspace{0.6cm}
\begin{center}
\textit{--- End of Module 3. Continued in Module 4: Problems, Submissions \& Code Execution Engine. ---}
\end{center}
\end{document}
